\documentclass[12pt]{article}

\usepackage[utf8]{inputenc}
\usepackage{amsmath, amssymb, amsfonts}
\usepackage{graphicx}
\usepackage[letterpaper, left=1in, right=1in, top=1in, bottom=1in]{geometry}
\usepackage{booktabs}
\usepackage{multirow}
\usepackage{float}
\usepackage{caption}
\usepackage{xcolor}
\usepackage{listings}
\usepackage{hyperref}
\usepackage[spanish]{babel}

\definecolor{codegray}{RGB}{245,245,245}
\definecolor{darkblue}{RGB}{0,45,90}

\lstset{
    backgroundcolor=\color{codegray},
    basicstyle=\ttfamily\small,
    breaklines=true,
    frame=single,
    columns=fullflexible,
    keepspaces=true
}

\title{
Guía de fundamentos para la configuración de redes\\[0.5em]
\large Administración de redes
}

\author{
Jorge Alberto Suarez Saldana\\
\small Matrícula: 355992\\
\small Profesor: Raymundo Antonio González Grimaldo
}

\date{\today}

\begin{document}

\maketitle

\tableofcontents
\newpage


\section{Introducción}

Una práctica de configuración de redes como la realizada en Cisco Packet
Tracer no consiste únicamente en memorizar comandos. Para poder construir
una red sin ayuda es necesario comprender cómo se relacionan los diferentes
elementos que la componen: dispositivos, interfaces, medios físicos,
direcciones IP, subredes, VLANs, protocolos de enlace, enrutamiento y
mecanismos de seguridad.

El objetivo de este documento es reunir los conocimientos fundamentales que
un ingeniero o estudiante de redes debería dominar para poder analizar y
configurar por sí mismo una topología similar a la utilizada en la práctica.

La idea fundamental es que, frente a una topología, se pueda responder
sistemáticamente a las siguientes preguntas:

\begin{enumerate}
    \item ¿Qué dispositivos existen y cuál es la función de cada uno?
    \item ¿Qué interfaces deben conectarse?
    \item ¿Qué tipo de cable se necesita?
    \item ¿Qué redes IP existen?
    \item ¿Qué dispositivos pertenecen a cada red?
    \item ¿Qué VLAN corresponde a cada dispositivo?
    \item ¿Qué enlaces deben ser \textit{access} y cuáles \textit{trunk}?
    \item ¿Dónde debe realizarse el enrutamiento?
    \item ¿Qué rutas necesita conocer cada router?
    \item ¿Qué medidas de seguridad deben configurarse?
    \item ¿Cómo se comprobará que la configuración funciona?
\end{enumerate}


\section{Modelo mental de una red}

Antes de comenzar a escribir comandos es conveniente entender una red como
una serie de capas.

Un dispositivo final genera información. Esa información se encapsula en
diferentes estructuras mientras atraviesa la red. Un switch toma decisiones
principalmente utilizando direcciones MAC, mientras que un router utiliza
direcciones IP para decidir hacia qué red debe enviar un paquete.

De forma simplificada:

\begin{center}
\textbf{Host} $\rightarrow$ \textbf{Switch} $\rightarrow$
\textbf{Router} $\rightarrow$ \textbf{Router} $\rightarrow$
\textbf{Switch} $\rightarrow$ \textbf{Host}
\end{center}

El switch permite la comunicación dentro de una red de capa 2, mientras que
el router permite comunicar diferentes redes de capa 3.

Esta distinción es fundamental. Si dos computadoras pertenecen a la misma
subred, normalmente pueden comunicarse mediante el switch sin necesitar
enrutamiento. Si pertenecen a subredes diferentes, necesitan un router o
algún dispositivo que realice una función equivalente.


\section{Dispositivos de red}

\subsection{PC, laptop y servidor}

Los dispositivos finales son los equipos que generan o reciben información.
En Packet Tracer pueden representarse mediante PCs, laptops y servidores.

Un host normalmente necesita como mínimo:

\begin{itemize}
    \item Dirección IP.
    \item Máscara de subred.
    \item Gateway predeterminado.
\end{itemize}

Por ejemplo:

\begin{lstlisting}
IP Address:       192.168.25.10
Subnet Mask:      255.255.255.0
Default Gateway:  192.168.25.1
\end{lstlisting}

La dirección IP identifica al dispositivo dentro de una red.

La máscara determina qué parte de la dirección corresponde a la red y qué
parte corresponde al host.

El gateway predeterminado es la dirección a la cual se envían los paquetes
destinados a otras redes.


\subsection{Switch}

Un switch es principalmente un dispositivo de capa 2. Su función principal
es conectar dispositivos dentro de una red local y reenviar tramas
utilizando direcciones MAC.

El switch aprende qué dirección MAC se encuentra asociada a cada puerto y
utiliza esa información para decidir por qué puerto enviar una trama.

En una práctica de VLANs, el switch además permite dividir lógicamente una
infraestructura física en diferentes redes de capa 2.


\subsection{Router}

Un router conecta diferentes redes IP.

Por ejemplo, si existen:

\[
192.168.25.0/24
\]

y

\[
192.168.30.0/24
\]

un router puede permitir la comunicación entre ellas.

Cada interfaz del router normalmente pertenece a una red diferente.

Por ejemplo:

\begin{lstlisting}
GigabitEthernet0/0
IP: 192.168.25.1

GigabitEthernet0/1
IP: 192.168.30.1
\end{lstlisting}

El router también mantiene una tabla de enrutamiento que le permite decidir
qué camino seguir para alcanzar una red determinada.


\section{Interfaces}

Las interfaces son los puntos de conexión de los dispositivos.

Algunos nombres comunes en routers Cisco son:

\begin{lstlisting}
GigabitEthernet0/0
GigabitEthernet0/1
Serial0/1/0
Serial0/1/1
\end{lstlisting}

En switches pueden aparecer interfaces como:

\begin{lstlisting}
FastEthernet0/1
FastEthernet0/2
GigabitEthernet0/1
\end{lstlisting}

Es importante leer correctamente la numeración de las interfaces. Conectar
un cable físicamente a un puerto no significa que ese puerto esté
correctamente configurado.

Para conocer el estado de las interfaces se utiliza:

\begin{lstlisting}
show ip interface brief
\end{lstlisting}

Una interfaz correctamente operativa normalmente debe mostrar:

\begin{lstlisting}
Status: up
Protocol: up
\end{lstlisting}

La combinación \texttt{up/up} indica que la interfaz está activa tanto a
nivel físico como a nivel de protocolo.


\section{Cables y conexiones}

Elegir correctamente el medio físico es una de las primeras decisiones que
deben tomarse al construir una topología.

\subsection{Cable de cobre directo}

El cable \textit{Copper Straight-Through} se utiliza tradicionalmente para
conectar dispositivos de diferente tipo, por ejemplo:

\begin{itemize}
    \item PC $\rightarrow$ Switch.
    \item Router $\rightarrow$ Switch.
\end{itemize}

En equipos modernos existen mecanismos como Auto-MDI/MDIX que permiten
detectar automáticamente el tipo de conexión, pero en prácticas de
Packet Tracer sigue siendo importante conocer la distinción.


\subsection{Cable cruzado}

El cable \textit{Copper Cross-Over} se utilizaba tradicionalmente para
conectar dispositivos del mismo tipo, por ejemplo:

\begin{itemize}
    \item Switch $\rightarrow$ Switch.
    \item Router $\rightarrow$ Router mediante Ethernet.
    \item PC $\rightarrow$ PC.
\end{itemize}

Nuevamente, los equipos modernos pueden realizar detección automática, pero
el concepto sigue siendo importante para comprender las conexiones Ethernet.


\subsection{Cable serial DCE y DTE}

Las conexiones seriales son diferentes de las conexiones Ethernet.

En una conexión serial existe un extremo DCE (\textit{Data Communications
Equipment}) y un extremo DTE (\textit{Data Terminal Equipment}).

El dispositivo DCE proporciona la señal de reloj necesaria para sincronizar
la comunicación.

Por esta razón, el comando:

\begin{lstlisting}
clock rate 64000
\end{lstlisting}

se configura únicamente en el extremo DCE.

Para verificar qué extremo es DCE puede utilizarse:

\begin{lstlisting}
show controllers serial 0/1/0
\end{lstlisting}

La salida permite identificar si el cable conectado es DCE o DTE.

Una regla práctica importante es:

\begin{center}
\textbf{DCE $\rightarrow$ configura clock rate}\\
\textbf{DTE $\rightarrow$ no configura clock rate}
\end{center}


\section{Direccionamiento IPv4}

Una dirección IPv4 tiene 32 bits y normalmente se representa mediante cuatro
octetos:

\[
192.168.25.10
\]

La máscara determina qué bits corresponden a la red.

Por ejemplo:

\[
192.168.25.10/24
\]

equivale a:

\[
255.255.255.0
\]

Una red /24 tiene:

\[
2^{32-24}=256
\]

direcciones totales.

De ellas, normalmente:

\[
256-2=254
\]

son direcciones utilizables para hosts.

La primera dirección identifica la red:

\[
192.168.25.0
\]

y la última corresponde al broadcast:

\[
192.168.25.255
\]


\section{Redes /30 para enlaces entre routers}

Los enlaces punto a punto entre routers no necesitan cientos de direcciones.
Por ello es común utilizar una red /30.

Por ejemplo:

\[
10.10.10.0/30
\]

tiene cuatro direcciones:

\begin{center}
\begin{tabular}{|c|c|}
\hline
\textbf{Dirección} & \textbf{Uso} \\
\hline
10.10.10.0 & Red \\
10.10.10.1 & Router 1 \\
10.10.10.2 & Router 2 \\
10.10.10.3 & Broadcast \\
\hline
\end{tabular}
\end{center}

Por eso una red /30 es muy conveniente para conectar dos routers.


\section{Cómo construir una tabla de direccionamiento}

Antes de configurar los dispositivos conviene elaborar una tabla de
direccionamiento.

Una tabla típica puede tener la siguiente estructura:

\begin{center}
\renewcommand{\arraystretch}{1.2}
\begin{tabular}{|l|l|l|l|l|}
\hline
\textbf{Dispositivo} & \textbf{Interfaz/VLAN} &
\textbf{Red} & \textbf{IP} & \textbf{Gateway} \\
\hline
R1 & G0/1 & 148.120.1.0/24 & 148.120.1.1 & -- \\
PC-A & LAN & 148.120.1.0/24 & 148.120.1.10 & 148.120.1.1 \\
\hline
R2 & VLAN 25 & 192.168.25.0/24 & 192.168.25.1 & -- \\
PC-D & VLAN 25 & 192.168.25.0/24 & 192.168.25.10 & 192.168.25.1 \\
\hline
R3 & VLAN 10 & 172.16.10.0/24 & 172.16.10.1 & -- \\
PC-E & VLAN 10 & 172.16.10.0/24 & 172.16.10.10 & 172.16.10.1 \\
\hline
\end{tabular}
\end{center}

La tabla debe elaborarse antes de configurar para evitar utilizar
direcciones duplicadas o asignar un gateway incorrecto.


\section{VLANs}

Una VLAN (\textit{Virtual Local Area Network}) permite dividir
lógicamente una infraestructura de red física.

Por ejemplo, un switch podría tener:

\begin{itemize}
    \item VLAN 10: Becarios.
    \item VLAN 20: Nativa.
    \item VLAN 25: Coordinadores.
    \item VLAN 30: Profesores.
\end{itemize}

Aunque los dispositivos estén conectados físicamente al mismo switch,
pertenecer a diferentes VLANs los coloca en diferentes dominios de
broadcast de capa 2.


\subsection{Crear una VLAN}

En un switch:

\begin{lstlisting}
Switch(config)#vlan 25
Switch(config-vlan)#name Coordinadores
Switch(config-vlan)#exit
\end{lstlisting}


\subsection{Puerto de acceso}

Un puerto conectado a un dispositivo final normalmente se configura como
\textit{access}:

\begin{lstlisting}
Switch(config)#interface fastEthernet 0/7
Switch(config-if)#switchport mode access
Switch(config-if)#switchport access vlan 25
\end{lstlisting}

De esta manera, el dispositivo conectado a ese puerto pertenece a la
VLAN 25.


\subsection{Puerto trunk}

Un enlace que necesita transportar varias VLANs debe utilizarse como
\textit{trunk}.

Por ejemplo:

\begin{lstlisting}
Switch(config)#interface gigabitEthernet 0/1
Switch(config-if)#switchport mode trunk
Switch(config-if)#switchport trunk native vlan 20
\end{lstlisting}

Los enlaces trunk son especialmente importantes cuando se conecta un switch
con otro switch o un switch con un router que utiliza
\textit{router-on-a-stick}.


\section{VLAN nativa}

En un enlace 802.1Q, una VLAN puede configurarse como VLAN nativa.

Por ejemplo:

\begin{lstlisting}
switchport trunk native vlan 20
\end{lstlisting}

En la práctica realizada se utilizó la VLAN 20 como VLAN nativa.

Es importante que ambos extremos del enlace trunk tengan configurada la
misma VLAN nativa para evitar inconsistencias en la comunicación.


\section{Router-on-a-stick}

Cuando se desea comunicar múltiples VLANs utilizando una sola interfaz física
del router, se utiliza una configuración conocida como
\textit{router-on-a-stick}.

La interfaz física se conecta mediante un enlace trunk al switch.

Por ejemplo:

\begin{lstlisting}
R2(config)#interface g0/1
R2(config-if)#no ip address
R2(config-if)#no shutdown
\end{lstlisting}

Después se crean subinterfaces:

\begin{lstlisting}
R2(config)#interface g0/1.25
R2(config-subif)#encapsulation dot1Q 25
R2(config-subif)#ip address 192.168.25.1 255.255.255.0
\end{lstlisting}

Y otra:

\begin{lstlisting}
R2(config)#interface g0/1.30
R2(config-subif)#encapsulation dot1Q 30
R2(config-subif)#ip address 192.168.30.1 255.255.255.0
\end{lstlisting}

Cada subinterfaz funciona como gateway de una VLAN diferente.

Por ejemplo:

\begin{center}
\begin{tabular}{|c|c|c|}
\hline
\textbf{VLAN} & \textbf{Red} & \textbf{Gateway} \\
\hline
25 & 192.168.25.0/24 & 192.168.25.1 \\
30 & 192.168.30.0/24 & 192.168.30.1 \\
\hline
\end{tabular}
\end{center}


\section{Enrutamiento}

El enrutamiento es el proceso mediante el cual un router determina por dónde
debe enviar un paquete para alcanzar una red determinada.

Un router necesita conocer las redes a las que está conectado directamente
y, si existen otras redes remotas, necesita aprender o recibir información
sobre cómo alcanzarlas.


\subsection{Redes directamente conectadas}

Cuando se configura una interfaz:

\begin{lstlisting}
interface g0/1
ip address 148.120.1.1 255.255.255.0
no shutdown
\end{lstlisting}

el router automáticamente conoce la red:

\[
148.120.1.0/24
\]

Esta aparecerá en la tabla de enrutamiento como una red conectada
directamente.


\subsection{Rutas estáticas}

Una ruta estática se configura manualmente.

La sintaxis general es:

\begin{lstlisting}
ip route RED MASCARA NEXT-HOP
\end{lstlisting}

Por ejemplo:

\begin{lstlisting}
R1(config)#ip route 192.168.25.0 255.255.255.0 10.10.10.2
\end{lstlisting}

Esto significa:

\begin{center}
``Para llegar a 192.168.25.0/24, envía el paquete hacia 10.10.10.2.''
\end{center}

Una ruta estática siempre debe analizarse desde la perspectiva del router
en el que se está configurando.


\section{Cómo diseñar las rutas estáticas}

Para construir rutas estáticas correctamente hay que pensar en ambos sentidos.

Supongamos:

\[
R1 \longleftrightarrow R2 \longleftrightarrow R3
\]

Si una computadora detrás de R1 quiere comunicarse con una computadora
detrás de R3, deben existir rutas para:

\begin{enumerate}
    \item R1 $\rightarrow$ R3.
    \item R2 $\rightarrow$ R3.
    \item R3 $\rightarrow$ R1.
\end{enumerate}

No basta con que el paquete llegue al destino. El destino también necesita
saber cómo devolver la respuesta.

Este es uno de los errores más comunes al configurar redes manualmente.


\section{Tabla de enrutamiento}

La tabla de enrutamiento puede consultarse mediante:

\begin{lstlisting}
show ip route
\end{lstlisting}

Algunas letras importantes son:

\begin{center}
\begin{tabular}{|c|l|}
\hline
\textbf{Código} & \textbf{Significado} \\
\hline
C & Connected, red conectada directamente \\
L & Local, dirección de una interfaz del router \\
S & Static, ruta estática \\
\hline
\end{tabular}
\end{center}

Por ejemplo:

\begin{lstlisting}
C    192.168.25.0/24 is directly connected
S    172.16.10.0/24 [1/0] via 10.10.20.2
\end{lstlisting}

La primera red está conectada directamente y la segunda se alcanza mediante
una ruta estática.


\section{Gateway predeterminado}

El gateway predeterminado es uno de los conceptos más importantes para
diagnosticar problemas de conectividad.

Si un host quiere comunicarse con otro dispositivo de la misma red, no
necesariamente necesita utilizar el gateway.

Pero si el destino pertenece a otra red, el host debe enviar el paquete al
gateway.

Por ejemplo:

\begin{lstlisting}
PC-D
IP:       192.168.25.10
Mask:     255.255.255.0
Gateway:  192.168.25.1
\end{lstlisting}

Si PC-D quiere comunicarse con:

\[
192.168.25.11
\]

el destino está en su misma red.

Pero si quiere comunicarse con:

\[
172.16.10.10
\]

debe enviar el paquete a:

\[
192.168.25.1
\]

Si el gateway está vacío o es incorrecto, la comunicación hacia otras redes
fallará aunque la configuración del switch y del router sea correcta.


\section{Seguridad básica de los dispositivos Cisco}

La seguridad de una red no comienza únicamente con protocolos avanzados.
También es necesario proteger el acceso administrativo a los dispositivos.


\subsection{Modo privilegiado}

El modo privilegiado se obtiene mediante:

\begin{lstlisting}
Router>enable
\end{lstlisting}

Para protegerlo puede utilizarse:

\begin{lstlisting}
Router(config)#enable secret class
\end{lstlisting}

El uso de \texttt{enable secret} es preferible a una contraseña configurada
simplemente con \texttt{enable password}.


\subsection{Contraseña de consola}

Puede protegerse el acceso por consola:

\begin{lstlisting}
Router(config)#line console 0
Router(config-line)#password cisco
Router(config-line)#login
\end{lstlisting}


\subsection{Cifrado de contraseñas}

Puede habilitarse:

\begin{lstlisting}
Router(config)#service password-encryption
\end{lstlisting}

Esto evita que ciertas contraseñas aparezcan directamente en texto plano
dentro de la configuración.


\subsection{Banner}

Un dispositivo también puede mostrar una advertencia al iniciar sesión:

\begin{lstlisting}
Router(config)#banner motd #ACCESO RESTRINGIDO - SOLO PERSONAL AUTORIZADO#
\end{lstlisting}

El banner tiene principalmente una función informativa y administrativa.


\section{Seguridad de puertos}

Un switch puede utilizar \textit{port-security} para limitar qué dispositivos
pueden utilizar un puerto de acceso.

Una configuración básica es:

\begin{lstlisting}
Switch(config)#interface fastEthernet 0/7
Switch(config-if)#switchport mode access
Switch(config-if)#switchport access vlan 25
Switch(config-if)#switchport port-security
Switch(config-if)#switchport port-security maximum 1
Switch(config-if)#switchport port-security violation shutdown
\end{lstlisting}

En este ejemplo se permite como máximo una dirección MAC en el puerto.

Si se detecta una violación, el puerto puede quedar deshabilitado.

La seguridad de puertos ayuda a impedir que cualquier dispositivo conectado
físicamente pueda utilizar libremente un puerto de la infraestructura.


\section{Modos de configuración de Cisco IOS}

Es fundamental distinguir los diferentes modos de la CLI.

\subsection{User EXEC}

\begin{lstlisting}
Router>
\end{lstlisting}

Permite realizar operaciones básicas.


\subsection{Privileged EXEC}

\begin{lstlisting}
Router#
\end{lstlisting}

Se obtiene con:

\begin{lstlisting}
enable
\end{lstlisting}

Aquí pueden ejecutarse comandos de diagnóstico como:

\begin{lstlisting}
show ip route
show running-config
show ip interface brief
\end{lstlisting}


\subsection{Global configuration}

\begin{lstlisting}
Router(config)#
\end{lstlisting}

Se obtiene mediante:

\begin{lstlisting}
configure terminal
\end{lstlisting}

Desde aquí se modifican configuraciones generales.


\subsection{Interface configuration}

\begin{lstlisting}
Router(config-if)#
\end{lstlisting}

Se obtiene mediante:

\begin{lstlisting}
interface g0/1
\end{lstlisting}

Aquí se configuran parámetros específicos de una interfaz.


\subsection{Subinterface configuration}

\begin{lstlisting}
Router(config-subif)#
\end{lstlisting}

Aparece al configurar una subinterfaz:

\begin{lstlisting}
interface g0/1.25
\end{lstlisting}


\section{Comandos fundamentales}

Los siguientes comandos deberían formar parte del repertorio básico de
cualquier estudiante de redes Cisco.

\subsection{Comandos de configuración}

\begin{lstlisting}
enable
configure terminal
hostname R1
no ip domain-lookup
interface g0/1
ip address 192.168.1.1 255.255.255.0
no shutdown
exit
end
\end{lstlisting}


\subsection{Comandos de diagnóstico}

\begin{lstlisting}
show ip interface brief
show running-config
show startup-config
show ip route
show vlan brief
show interfaces trunk
\end{lstlisting}


\subsection{Comandos de conectividad}

Desde un host:

\begin{lstlisting}
ping 192.168.25.1
ping 192.168.30.10
\end{lstlisting}

Desde un router:

\begin{lstlisting}
ping 202.20.3.10
\end{lstlisting}

También puede utilizarse:

\begin{lstlisting}
traceroute 202.20.3.10
\end{lstlisting}

para observar por qué dispositivos atraviesa el paquete.


\subsection{Guardar la configuración}

La configuración activa se encuentra en la memoria RAM y puede perderse
si el dispositivo se reinicia.

Para guardarla:

\begin{lstlisting}
copy running-config startup-config
\end{lstlisting}

La idea fundamental es:

\begin{center}
\texttt{running-config} = configuración activa

\texttt{startup-config} = configuración guardada para el arranque
\end{center}


\section{Diagnóstico sistemático de problemas}

Cuando un ping falla, no conviene comenzar cambiando rutas al azar.

Un ingeniero de redes debe diagnosticar de abajo hacia arriba.


\subsection{Paso 1: revisar el cableado}

Primero debe verificarse:

\begin{itemize}
    \item ¿El cable conecta los puertos correctos?
    \item ¿El tipo de cable es apropiado?
    \item ¿Las interfaces están físicamente activas?
\end{itemize}


\subsection{Paso 2: revisar interfaces}

En routers:

\begin{lstlisting}
show ip interface brief
\end{lstlisting}

Se debe buscar especialmente:

\begin{lstlisting}
up    up
\end{lstlisting}

Si aparece:

\begin{lstlisting}
administratively down
\end{lstlisting}

normalmente falta:

\begin{lstlisting}
no shutdown
\end{lstlisting}


\subsection{Paso 3: comprobar la configuración IP}

Verificar:

\begin{itemize}
    \item Dirección IP.
    \item Máscara.
    \item Gateway.
\end{itemize}

Un error muy pequeño en cualquiera de estos valores puede impedir la
comunicación.


\subsection{Paso 4: comprobar las VLANs}

En el switch:

\begin{lstlisting}
show vlan brief
\end{lstlisting}

Esto permite verificar qué VLANs existen y qué puertos pertenecen a cada
una.


\subsection{Paso 5: comprobar los trunks}

Utilizar:

\begin{lstlisting}
show interfaces trunk
\end{lstlisting}

Se debe comprobar que:

\begin{itemize}
    \item El puerto realmente funciona como trunk.
    \item Las VLAN necesarias están permitidas.
    \item La VLAN nativa es consistente.
\end{itemize}


\subsection{Paso 6: comprobar el gateway}

Desde el host se puede hacer primero ping al gateway.

Por ejemplo:

\begin{lstlisting}
ping 192.168.25.1
\end{lstlisting}

Si este ping falla, todavía no tiene sentido investigar rutas entre routers:
el problema se encuentra probablemente entre el host y su gateway.


\subsection{Paso 7: comprobar el enrutamiento}

En cada router:

\begin{lstlisting}
show ip route
\end{lstlisting}

Hay que preguntarse:

\begin{center}
\textbf{¿El router conoce la red destino?}
\end{center}

Y posteriormente:

\begin{center}
\textbf{¿El router destino sabe cómo regresar?}
\end{center}


\section{Metodología para configurar una topología desde cero}

Una metodología ordenada reduce considerablemente los errores.

\subsection{Paso 1: estudiar la topología}

Antes de introducir comandos se deben identificar:

\begin{itemize}
    \item Routers.
    \item Switches.
    \item Hosts.
    \item Servidores.
    \item Interfaces.
    \item Enlaces.
\end{itemize}


\subsection{Paso 2: identificar las redes}

Cada segmento de capa 3 debe tener su propia red.

Por ejemplo:

\begin{center}
\begin{tabular}{|c|c|}
\hline
\textbf{Segmento} & \textbf{Red} \\
\hline
LAN R1 & 148.120.1.0/24 \\
R1--R2 & 10.10.10.0/30 \\
LAN R2 & 148.120.2.0/24 \\
R2--R3 & 10.10.20.0/30 \\
VLAN 25 & 192.168.25.0/24 \\
VLAN 30 & 192.168.30.0/24 \\
VLAN 10 & 172.16.10.0/24 \\
Servidor & 202.20.3.0/24 \\
\hline
\end{tabular}
\end{center}


\subsection{Paso 3: construir la tabla de direccionamiento}

Nunca debería ser necesario ``adivinar'' una dirección IP mientras se
configura un dispositivo.

La tabla debe contener todas las direcciones antes de comenzar.


\subsection{Paso 4: configurar los dispositivos}

Una posible secuencia es:

\begin{enumerate}
    \item Configuración básica de routers.
    \item Configuración básica de switches.
    \item Direcciones IP de interfaces.
    \item VLANs.
    \item Puertos de acceso.
    \item Trunks.
    \item Subinterfaces.
    \item Enlaces seriales.
    \item Rutas estáticas.
    \item Seguridad.
\end{enumerate}


\subsection{Paso 5: configurar los hosts}

Cada host debe recibir:

\begin{lstlisting}
IP Address
Subnet Mask
Default Gateway
\end{lstlisting}

El gateway debe corresponder con la interfaz del router perteneciente a la
misma red.


\subsection{Paso 6: probar progresivamente}

No es recomendable probar únicamente un ping entre dos extremos de una
topología compleja.

Primero:

\[
Host \rightarrow Gateway
\]

Después:

\[
Router \rightarrow Router
\]

Después:

\[
Host \rightarrow Host
\]

Finalmente se prueban comunicaciones entre redes alejadas.


\section{Ejemplo de razonamiento completo}

Supongamos que PC-C tiene:

\begin{lstlisting}
IP:       192.168.30.11
Mask:     255.255.255.0
Gateway:  192.168.30.1
\end{lstlisting}

y PC-E:

\begin{lstlisting}
IP:       172.16.10.10
Mask:     255.255.255.0
Gateway:  172.16.10.1
\end{lstlisting}

PC-C y PC-E pertenecen a redes diferentes.

Por lo tanto:

\[
192.168.30.0/24
\neq
172.16.10.0/24
\]

PC-C debe enviar el paquete a su gateway:

\[
192.168.30.1
\]

El router R2 debe saber cómo llegar a:

\[
172.16.10.0/24
\]

Por ejemplo:

\begin{lstlisting}
ip route 172.16.10.0 255.255.255.0 10.10.20.2
\end{lstlisting}

Después, R3 debe recibir el paquete y entregarlo mediante su interfaz
correspondiente a la VLAN 10.

Para que la respuesta funcione, R3 también debe conocer cómo regresar a
192.168.30.0/24.

Este ejemplo muestra por qué la conectividad de una red depende de toda la
cadena y no únicamente de la configuración del dispositivo que contiene
al host.


\section{Cómo agregar una nueva VLAN}

Supongamos que se solicita:

\[
\text{VLAN 40 -- Alumnos}
\]

con:

\[
192.168.40.0/24
\]

El procedimiento general sería:

\begin{enumerate}
    \item Crear la VLAN 40 en los switches necesarios.
    \item Asignar los puertos de los alumnos a la VLAN 40.
    \item Permitir la VLAN 40 por los enlaces trunk necesarios.
    \item Crear una subinterfaz para VLAN 40 en el router.
    \item Asignarle un gateway, por ejemplo 192.168.40.1.
    \item Configurar los hosts de la nueva VLAN.
    \item Agregar rutas estáticas en otros routers si necesitan alcanzar
    la nueva red.
    \item Probar conectividad.
\end{enumerate}

La subinterfaz podría ser:

\begin{lstlisting}
interface g0/1.40
encapsulation dot1Q 40
ip address 192.168.40.1 255.255.255.0
\end{lstlisting}

Este procedimiento es una aplicación directa de los conceptos de VLAN,
trunk, router-on-a-stick y enrutamiento.


\section{Tabla de comandos esenciales}

\begin{center}
\renewcommand{\arraystretch}{1.3}
\begin{tabular}{|l|p{8cm}|}
\hline
\textbf{Comando} & \textbf{Función} \\
\hline
\texttt{enable} & Entrar al modo privilegiado. \\
\hline
\texttt{configure terminal} & Entrar al modo de configuración global. \\
\hline
\texttt{hostname} & Cambiar el nombre del dispositivo. \\
\hline
\texttt{interface} & Entrar a una interfaz. \\
\hline
\texttt{ip address} & Configurar una dirección IP. \\
\hline
\texttt{no shutdown} & Activar una interfaz. \\
\hline
\texttt{shutdown} & Desactivar una interfaz. \\
\hline
\texttt{vlan} & Crear o modificar una VLAN. \\
\hline
\texttt{switchport mode access} & Configurar un puerto como acceso. \\
\hline
\texttt{switchport mode trunk} & Configurar un puerto como trunk. \\
\hline
\texttt{switchport access vlan} & Asignar un puerto de acceso a una VLAN. \\
\hline
\texttt{encapsulation dot1Q} & Asociar una subinterfaz con una VLAN. \\
\hline
\texttt{ip route} & Crear una ruta estática. \\
\hline
\texttt{show ip interface brief} & Mostrar estado resumido de interfaces. \\
\hline
\texttt{show ip route} & Mostrar tabla de enrutamiento. \\
\hline
\texttt{show running-config} & Mostrar configuración activa. \\
\hline
\texttt{show vlan brief} & Mostrar VLANs y puertos asociados. \\
\hline
\texttt{show interfaces trunk} & Mostrar información de enlaces trunk. \\
\hline
\texttt{ping} & Comprobar conectividad IP. \\
\hline
\texttt{traceroute} & Mostrar el camino hacia un destino. \\
\hline
\texttt{copy running-config startup-config} & Guardar la configuración. \\
\hline
\end{tabular}
\end{center}


\section{Lista de comprobación}

Antes de considerar terminada una práctica de este tipo, conviene realizar
la siguiente comprobación.

\subsection*{Capa física}

\begin{itemize}
    \item[$\square$] Los cables están conectados a los puertos correctos.
    \item[$\square$] Los enlaces seriales tienen un extremo DCE.
    \item[$\square$] El extremo DCE tiene configurado \texttt{clock rate}.
    \item[$\square$] Las interfaces necesarias están en estado \texttt{up/up}.
\end{itemize}

\subsection*{Direccionamiento}

\begin{itemize}
    \item[$\square$] Cada red tiene una dirección diferente.
    \item[$\square$] No existen direcciones IP duplicadas.
    \item[$\square$] Las máscaras son correctas.
    \item[$\square$] Cada host tiene un gateway correcto.
\end{itemize}

\subsection*{VLANs}

\begin{itemize}
    \item[$\square$] Todas las VLANs necesarias fueron creadas.
    \item[$\square$] Los puertos de acceso pertenecen a la VLAN correcta.
    \item[$\square$] Los enlaces necesarios son trunk.
    \item[$\square$] La VLAN nativa coincide en ambos extremos.
\end{itemize}

\subsection*{Enrutamiento}

\begin{itemize}
    \item[$\square$] Cada router conoce sus redes directamente conectadas.
    \item[$\square$] Existen rutas hacia las redes remotas.
    \item[$\square$] Las rutas de regreso también existen.
    \item[$\square$] Los siguientes saltos son alcanzables.
\end{itemize}

\subsection*{Seguridad}

\begin{itemize}
    \item[$\square$] Existe una contraseña de modo privilegiado.
    \item[$\square$] El acceso por consola está protegido.
    \item[$\square$] Las contraseñas correspondientes están protegidas.
    \item[$\square$] Los puertos de acceso tienen \textit{port-security} cuando
    sea necesario.
\end{itemize}

\subsection*{Pruebas}

\begin{itemize}
    \item[$\square$] Cada host puede alcanzar su gateway.
    \item[$\square$] Los routers pueden comunicarse entre sí.
    \item[$\square$] Los hosts pueden comunicarse dentro de su VLAN.
    \item[$\square$] Los hosts pueden comunicarse entre diferentes VLANs.
    \item[$\square$] Se puede alcanzar el servidor desde las diferentes redes.
\end{itemize}


\section{Conclusión}

Configurar una red correctamente requiere mucho más que conocer la sintaxis
de Cisco IOS. El verdadero objetivo es desarrollar un modelo mental que
permita interpretar una topología y convertirla en una configuración
coherente.

Los conceptos fundamentales están relacionados entre sí. Los cables
permiten construir la conectividad física; los switches permiten transportar
tramas y separar dominios mediante VLANs; los routers conectan diferentes
redes; las direcciones IP identifican los dispositivos dentro de esas redes;
y las tablas de enrutamiento permiten determinar hacia dónde deben enviarse
los paquetes.

De la misma manera, conceptos como \textit{trunk}, VLAN nativa y
\textit{router-on-a-stick} no deben aprenderse como comandos independientes.
Todos describen diferentes partes del mismo problema: transportar y enrutar
tráfico perteneciente a múltiples redes sobre una infraestructura física
compartida.

Una buena metodología consiste en comenzar siempre por la topología,
identificar las redes, elaborar una tabla de direccionamiento y después
configurar progresivamente los dispositivos. Finalmente, la conectividad
debe comprobarse desde los elementos más cercanos hasta los más lejanos.

Cuando un ping falla, el objetivo no debe ser simplemente ``hacer que
funcione'', sino determinar en qué punto de la cadena se interrumpe la
comunicación. Para ello se deben revisar sucesivamente la capa física, las
interfaces, las VLANs, los gateways, las tablas de enrutamiento y las rutas
de regreso.

Dominar esta forma de pensar permite pasar de seguir instrucciones
específicas de una práctica a ser capaz de analizar y construir una red por
cuenta propia.

\end{document}